\chapter{两翼算符与开放网络的深入推导}\label{app:wing-operators-networks}

本附录为正文第~4.6~节（同核异能态四大物理机制）、第~4.7~节（弱衰变多极算符与双贝塔微观矩阵元）以及第~4.8~节（开放系统与环境布居动力学）提供自洽、严格的数学算符推导、代数相消证明、马尔可夫吸收网络分析以及开放主方程本征模分解。
为了保持正文物理图像的紧凑性与连贯性，涉及准自旋群代数 Wigner-Eckart 约化、高阶禁戒形状因子、双贝塔衰变短程手征接触项、多步相继缓发中子吸收链以及非平衡双能级动力学的完整推演统一在此收录。

\section{单 $j$ 壳层准自旋 $\mathrm{SU}(2)$ 代数与同资数相消定理}\label{sec:appD-seniority-quasispin}

\subsection{准自旋算符与李代数换位关系}

考虑由单粒子角动量为 $j$ 的单壳层构成的费米子体系。壳层的单粒子态对简并度定义为：
\begin{equation}\label{eq:appD-omega-def}
\Omega_j \equiv j + \frac{1}{2} = \frac{2j + 1}{2}.
\end{equation}
为了描述核子之间的短程配对相互作用，定义单壳层双核子配对产生算符与湮灭算符为~\cite{book-jain-nuclear-isomers, book-suhonen-from-nucleons-to-nucleus}：
\begin{align}
\hat{S}_+ &\equiv \frac{1}{2} \sum_{m=-j}^j (-1)^{j-m} a_{jm}^\dagger a_{j,-m}^\dagger, \label{eq:appD-splus-def} \\
\hat{S}_- &\equiv (\hat{S}_+)^\dagger = \frac{1}{2} \sum_{m=-j}^j (-1)^{j-m} a_{j,-m} a_{jm}. \label{eq:appD-sminus-def}
\end{align}
对应的中性算符由总粒子数算符 $\hat{N} = \sum_m a_{jm}^\dagger a_{jm}$ 定义为：
\begin{equation}\label{eq:appD-szero-def}
\hat{S}_0 \equiv \frac{1}{2} \left( \hat{N} - \Omega_j \right).
\end{equation}
利用标准费米子反对易关系 $\{a_{jm}, a_{jm'}^\dagger\} = \delta_{mm'}$，可直接计算算符之间的对易子：
\begin{equation}\label{eq:appD-quasispin-commutation}
[\hat{S}_+, \hat{S}_-] = 2 \hat{S}_0, \qquad [\hat{S}_0, \hat{S}_\pm] = \pm \hat{S}_\pm.
\end{equation}
式~\eqref{eq:appD-quasispin-commutation} 严格构成了标准李群 $\mathrm{SU}(2)$ 的李代数。因此，多体配对态完全可以按照准自旋角动量理论进行不可约分类。

\subsection{量子数分类、阶梯结构与 Seniority 资数}

设准自旋 Casimir 算符为 $\hat{\mathbf{S}}^2 = \hat{S}_+ \hat{S}_- + \hat{S}_0^2 - \hat{S}_0$。本征态用准自旋量子数 $S$ 及其第三分量投影 $S_0$ 标记为 $|S, S_0\rangle$。
定义 Seniority（资数）$v$ 为多体态中未配对（未被 $\hat{S}_+$ 束缚成 $J=0$ 对）的价核子数目。由于资数为 $v$ 的基态 $|j^v, v, I\rangle$ 无法再被对湮灭算符破坏：
\begin{equation}\label{eq:appD-seniority-ground-cond}
\hat{S}_- |j^v, v, I\rangle = 0,
\end{equation}
该态对应于准自旋多重态的最底端态（最低权重态），其粒子数 $n = v$。由式~\eqref{eq:appD-szero-def}，其投影为：
\begin{equation}\label{eq:appD-s0-lowest}
S_0^{\min} = \frac{1}{2}(v - \Omega_j) = -S \implies S = \frac{1}{2}(\Omega_j - v).
\end{equation}
通过作用对产生算符 $\hat{S}_+$，粒子数按步长 2 递增，生成包含 $(n-v)/2$ 个配对核子的激发多体阶梯态：
\begin{equation}\label{eq:appD-quasispin-ladder}
n = v, v + 2, v + 4, \dots, 2\Omega_j - v,
\end{equation}
第三分量取值为 $S_0 = \frac{1}{2}(n - \Omega_j) \in [-S, +S]$。

\subsection{单体电多极算符的准自旋张量性质与相消证明}

任意单体无自旋翻转的电磁多极算符在单 $j$ 壳层中展开为：
\begin{equation}\label{eq:appD-one-body-em-op}
\hat{T}^{(\lambda)}_\mu = \sum_{m, m'} \langle j m | t^{(\lambda)}_\mu | j m' \rangle a_{jm}^\dagger a_{jm'}.
\end{equation}
在准自旋空间中，粒子产生算符与时间反演湮灭算符组成准自旋二重态算符。通过反对易代数分解，偶多极阶数（$\lambda = 2, 4, \dots$ 为偶数，如核心的电四极算符 $E2$）的单体算符在准自旋空间严格变换为\emph{准自旋矢量}（张量秩 $\kappa = 1$ 的算符，其第三分量为 $\kappa_0 = 0$）：
\begin{equation}\label{eq:appD-e2-quasispin-vector}
\mathcal{M}(E2) \equiv \hat{T}^{(\kappa=1, \kappa_0=0)}.
\end{equation}
根据准自旋空间中的 Wigner-Eckart 定理，任意两个同资数（$\Delta v = 0 \implies \Delta S = 0$）多体态之间的电四极矩阵元，其粒子数依赖性完全由准自旋空间中的 Clebsch-Gordan 系数决定：
\begin{equation}\label{eq:appD-wigner-eckart-quasispin}
\langle j^n, v, I_f \| \mathcal{M}(E2) \| j^n, v, I_i \rangle = \frac{\langle S, S_0, 1, 0 \mid S, S_0 \rangle}{\langle S, S, 1, 0 \mid S, S \rangle} \langle j^v, v, I_f \| \mathcal{M}(E2) \| j^v, v, I_i \rangle.
\end{equation}
利用准自旋矢量 Clebsch-Gordan 系数的显式代数解析表达式：
\begin{equation}\label{eq:appD-cg-analytical}
\langle S, S_0, 1, 0 \mid S, S_0 \rangle = \frac{S_0}{\sqrt{S(S+1)}}, \qquad \langle S, S, 1, 0 \mid S, S \rangle = \frac{S}{\sqrt{S(S+1)}}.
\end{equation}
代入式~\eqref{eq:appD-s0-lowest} 与 $S_0 = \frac{1}{2}(n - \Omega_j)$，两项之比精确化简为：
\begin{equation}\label{eq:appD-cg-ratio}
\frac{\langle S, S_0, 1, 0 \mid S, S_0 \rangle}{\langle S, S, 1, 0 \mid S, S \rangle} = \frac{S_0}{S} = \frac{\frac{1}{2}(n - \Omega_j)}{\frac{1}{2}(\Omega_j - v)} = \frac{n - \Omega_j}{\Omega_j - v}.
\end{equation}
因此，同资数电四极矩阵元与基准态矩阵元之间满足普适的单调线性缩放定律：
\begin{equation}\label{eq:appD-seniority-e2-scaling}
\langle j^n, v, I_f \| \mathcal{M}(E2) \| j^n, v, I_i \rangle = \frac{\Omega_j - n}{\Omega_j - v} \langle j^v, v, I_f \| \mathcal{M}(E2) \| j^v, v, I_i \rangle.
\end{equation}

\subsection{半满相消定理与物理边界}

式~\eqref{eq:appD-seniority-e2-scaling} 导出了核结构中著名的\textbf{单壳层同资数电四极跃迁相消定理}：
当壳层中的核子数恰好等于对简并度时，即系统处于\textbf{半满壳层（Half-Filled Shell, $n = \Omega_j$）}：
\begin{equation}\label{eq:appD-cancellation-theorem}
\Omega_j - n = 0 \implies \langle j^{\Omega_j}, v, I_f \| \mathcal{M}(E2) \| j^{\Omega_j}, v, I_i \rangle \equiv 0.
\end{equation}
在半满壳层处，电四极跃迁几率 $B(E2; v \to v) \propto |\langle \mathcal{M}(E2) \rangle|^2 \equiv 0$ 精确消失！

在将该定理应用于实验分析时，必须严格遵守以下物理边界条件：
\begin{enumerate}
\item 资数奇异点处理（$v = \Omega_j$）：当 $v = \Omega_j$ 时，准自旋 $S = 0$，体系为准自旋单态，其粒子数阶梯仅允许单一取值 $n = \Omega_j$。式~\eqref{eq:appD-seniority-e2-scaling} 分母为零仅为表示式奇异，由于准自旋标量态之间跃迁张量三角法则 $|0 - 0| \le 1 \le 0$ 不满足，该单态内同资数跃迁本征为零，不存在真实发散；
\item 跃迁相消不等于总衰变宽度为零：相消定理仅预言了同资数 $\Delta v = 0$ 的电四极跃迁被压低；核态的总衰变依然可以通过磁偶极（$M1$）、高阶电十六极（$E4$）、内转换电子过程或微弱的跨资数（$\Delta v = 2$）混合衰变。半满相消造就的是寿命显著拉长（达纳秒至微秒量级）的 Seniority 同核异能态（如 $^{210}\mathrm{Po}$ 的 $8^+$ 态），绝非永不衰变的真稳定态；
\item 多壳层混合的破缺：若价核子处于多个相互纠缠的非简并壳层（如 $1f_{5/2}, 2p_{3/2}, 1g_{9/2}$），广义 Seniority 必须通过 Bogoliubov 准粒子平均场或大尺度壳模型修正，纯半满相消点可能发生显著偏移。
\end{enumerate}

\section{$K$ 禁戒度、Weisskopf 阻碍与非正交投影展开}\label{sec:appD-k-forbidden}

\subsection{轴对称转子模型 $K$ 禁戒度与 Weisskopf 阻碍因子}

在轴对称形变核中，核总角动量沿对称轴的投影量子数记为 $K$。对于纯辐射多极阶数为 $\lambda$ 的电磁跃迁，初末态投影差必须满足角动量矢量耦合选择定则：
\begin{equation}\label{eq:appD-k-selection-rule}
\Delta K \equiv |K_i - K_f| \le \lambda.
\end{equation}
当两个高 $K$ 组态之间发生退激、且其投影差超出辐射多极阶数时，定义\textbf{$K$ 禁戒度}为：
\begin{equation}\label{eq:appD-k-forbiddenness-nu}
\nu \equiv \Delta K - \lambda = |K_i - K_f| - \lambda > 0.
\end{equation}

为了量化跃迁受阻碍的程度，实验上将测得的偏 $\gamma$ 跃迁半衰期 $T_{1/2}^\gamma$ 与 Weisskopf 单粒子单极半衰期估计值 $T_{1/2}^W(\lambda)$ 进行对比，定义实验 Weisskopf 阻碍因子（Hindrance Factor）为：
\begin{equation}\label{eq:appD-hindrance-factor}
F_W \equiv \frac{T_{1/2}^\gamma}{T_{1/2}^W(\lambda)} = \frac{B_W(\lambda)}{B(\lambda; I_i \to I_f)}.
\end{equation}
其中实验部分光子半衰期由测得的核态平均寿命 $\tau$、总分支比 $b_\gamma$ 及全内转换系数 $\alpha_{\mathrm{tot}}$ 严格确定：
\begin{equation}\label{eq:appD-t12-gamma}
T_{1/2}^\gamma = \tau \ln 2 \cdot \frac{1 + \alpha_{\mathrm{tot}}}{b_\gamma}.
\end{equation}

\subsection{平均折合阻碍度 $f_\nu$ 的物理性质与局限}

为了在不同禁戒阶数 $\nu$ 的异能态之间建立可比标度，广泛采用经验的折合阻碍度（Reduced Hindrance per Degree of Forbiddenness）：
\begin{equation}\label{eq:appD-fnu-def}
f_\nu \equiv (F_W)^{1/\nu} = \left( \frac{T_{1/2}^\gamma}{T_{1/2}^W} \right)^{1/\nu}.
\end{equation}
在理想的刚性轴对称转子模型中，每跨越一个禁戒单位，微观波函数重叠受到 Corioils 相互作用或四极扰动的微扰衰减，使得 $f_\nu$ 在经验统计上保持在 $10$ 至 $100$ 范围（Jain et al. 2015 统计规范~\cite{book-jain-nuclear-isomers}）。
然而，必须在理论上严格界定：
\begin{enumerate}
\item $f_\nu$ 是实验宏观跃迁率的几何平均表征，\emph{绝对不是微观量子电动力学的第一性原理守恒常数}；
\item $f_\nu$ 的急剧崩塌是指示非轴对称形变的敏感探针：当原子核具备显著的三轴形变（$\gamma \ne 0^\circ$）时，非轴对称势阱引入强烈的 $\Delta K = \pm 2$ 组态相干混合，波函数中的微小杂质分量开通了低禁戒度的旁路跃迁通道，导致表观 $f_\nu$ 急剧下降至小于 5 甚至接近 1。因此，不能仅凭孤立谱线的较小辐射几率断言其形变性质。
\end{enumerate}

\subsection{非正交基下的两端态投影展开与范数核接口}

当核多体波函数由三维角动量投影自约束内禀 Hartre-Fock-Bogoliubov (HFB) 组态生成时（见附录 B），初态 $|\Psi_i\rangle$ 与末态 $|\Psi_f\rangle$ 的内禀轴对称对称性可能不同。跃迁算符在投影非正交态之间的约化矩阵元必须计入重叠范数核：
\begin{equation}\label{eq:appD-projected-nme}
\langle \Psi_f^{I_f K_f} \| \hat{\mathcal{M}}(\lambda) \| \Psi_i^{I_i K_i} \rangle = \frac{\int \mathrm{d}\Omega \, D^{I_f *}_{M_f K_f}(\Omega) \langle \Phi_f | \hat{\mathcal{M}}(\lambda) \hat{R}(\Omega) | \Phi_i \rangle}{\sqrt{\mathcal{N}_{I_f K_f} \mathcal{N}_{I_i K_i}}},
\end{equation}
式中 $\mathcal{N}_{IK}$ 为附录 B 式~\eqref{eq:appB-haar-measure} 定义的三维角动量投影范数核积分。投影算符的非正交性自然包含了内禀态之间的几何错位与组态重叠，消除了简单绝热近似在小形变下的代数发散。

\section{弱相互作用多极算符分解与双贝塔短程接触项}\label{sec:appD-weak-operators}

\subsection{闭合单体算符和定则与相空间积分率的解耦}

在弱电同位旋代数中，裸单体 Gamow-Teller 算符定义为 $\hat{\boldsymbol{\beta}}^\mp_{GT} = \sum_k \boldsymbol{\sigma}(k) t_\mp(k)$。利用同位旋算符的标准对易关系 $[t_-, t_+] = 2 t_3$，在无截断的全状态空间闭合条件下，Ikeda 差和定则严格导出为空间-同位旋几何恒等式（第~2.7~节式~\eqref{eq:2-7-ikeda-result}）：
\begin{equation}\label{eq:appD-ikeda-sum-rule}
S_- - S_+ \equiv \sum_{f_-} |\langle f_- \| \hat{\boldsymbol{\beta}}^-_{GT} \| i \rangle|^2 - \sum_{f_+} |\langle f_+ \| \hat{\boldsymbol{\beta}}^+_{GT} \| i \rangle|^2 = 3(N - Z).
\end{equation}
式~\eqref{eq:appD-ikeda-sum-rule} 完全独立于外在弱相互作用耦合常数 $g_A^2$，亦未包含任何轻子运动学因子的加权。

相反，实验上测得的物理衰变常数 $\lambda = \ln 2 / T_{1/2}$ 是跃迁强度在有限衰变能释放窗 $Q_\beta$ 内与微分轻子相空间加权积分的宏观结果：
\begin{equation}\label{eq:appD-lambda-phase-space}
\lambda = \frac{(G_F^{(E)})^2 |V_{ud}|^2 (m_e c^2)^5}{2\pi^3 \hbar} \int_1^{W_0} F(Z, W) p W (W_0 - W)^2 C(W) \, \mathrm{d}W.
\end{equation}
因此，有限价空间截断不仅造成高能巨共振强度的遗漏，有限的动能释放窗 $Q_\beta$ 更是使得绝大多数核素仅能观测到极小一部分允许或禁戒跃迁强度。不能将低能衰变强度的不足简单解释为基本弱流耦合常数的微观改变。

\subsection{一阶禁戒衰变算符的四维张量秩分解}

当允许 Gamow-Teller 与 Fermi 跃迁被宇称选择定则严格禁戒（初末态宇称反转 $\pi_i \pi_f = -1$）时，弱衰变主要由一阶禁戒（First-Forbidden, FF）通道驱动。
根据非相对论多极展开，一阶禁戒算符由四个具有不同空间对称性与张量秩的微观算符构成~\cite{doctoral-r1-g34-zhi-2013-prc-shell-model-half-lives}：
\begin{enumerate}
\item 秩为 0 的类时轴荷算符：
\begin{equation}\label{eq:appD-ff-rank0}
\hat{\mathcal{O}}_0 \propto \gamma_5 \to \frac{\boldsymbol{\sigma} \cdot \mathbf{p}}{m_N};
\end{equation}
\item 秩为 1 的矢量算符与轴矢量算符：
\begin{equation}\label{eq:appD-ff-rank1}
\hat{\mathcal{O}}_1^V \propto \boldsymbol{\alpha} \to \frac{\mathbf{p}}{m_N}, \qquad \hat{\mathcal{O}}_1^A \propto \boldsymbol{\sigma} \times \mathbf{r}, \quad \mathbf{r};
\end{equation}
\item 秩为 2 的空间张量算符：
\begin{equation}\label{eq:appD-ff-rank2}
\hat{\mathcal{O}}_2 \propto [\boldsymbol{\sigma} \otimes \mathbf{r}]^{(2)}.
\end{equation}
\end{enumerate}
各张量分量通过出射轻子的大分量与小分量狄拉克波函数干涉，在式~\eqref{eq:appD-lambda-phase-space} 中合成能量依赖的非平坦谱形状因子：
\begin{equation}\label{eq:appD-ff-shape-factor}
C(W) = k_0 + k_1 W + k_{-1} W^{-1} + k_2 W^2,
\end{equation}
各系数 $k_i$ 包含不同秩微观核矩阵元之间的相干干涉项。在 $N=126$ 等待点核区，由于中子跨壳跃迁（如 $\nu 1i_{13/2} \to \pi 1h_{9/2}$）受泡利阻塞压低，一阶禁戒跃迁显著增强，成为主导快中子俘获过程寿命的关键机制。

\subsection{双贝塔衰变矩阵元与手征有效场论短程接触算符}

在轻 Majorana 中微子交换机制下，$0\nu\beta\beta$ 的理论半衰期由式~\eqref{eq:4-7-0nubb-halflife} 给出。由于两个中微子在核内两个核子之间虚交换，传播子积分跨越极其宽广的虚拟动量谱：
\begin{equation}\label{eq:appD-0nubb-virtual-momentum}
|\mathbf{q}| \sim 100\text{--}200\,\mathrm{MeV}/c,
\end{equation}
探测的是 $r \sim 1$–$2\,\text{fm}$ 的短程核结构。

传统的核矩阵元通常分解为长程类库仑势积分：
\begin{equation}\label{eq:appD-0nubb-long-range}
M_{0\nu}^{\mathrm{long}} = M_{GT} - \left(\frac{g_V}{g_A}\right)^2 M_F + M_T.
\end{equation}
然而，Cirigliano 等人（2018）基于手征有效场论（Chiral EFT）的第一性原理证明~\cite{cirigliano-2018-leading-short-range-double-beta}：当在非微扰散射理论下求解核子间的双核子 $^1S_0$ 状态时，单独的长程单中微子交换势 $V_\nu^{\mathrm{long}} \sim 1/r_{12}$ 将在薛定谔方程中产生未被吸收的紫外发散。为了使散射振幅与核矩阵元满足微观重整化群不变性，必须在领头阶（Leading Order, LO）显式引入一个 $\Delta L = 2$ 的两核子短程接触相互作用势：
\begin{equation}\label{eq:appD-contact-potential}
V_\nu^{\mathrm{short}}(\mathbf{r}_{12}) = -2 g_\nu^{NN} \tau_{1,+} \tau_{2,+} \delta^{(3)}(\mathbf{r}_{12}),
\end{equation}
其中低能耦合常数 $g_\nu^{NN}$ 必须通过同位旋破坏电磁质量分裂或格点 QCD 匹配进行独立标定。
必须严肃明确两项科学边界：
\begin{enumerate}
\item 该短程接触算符是理论自洽重整化的必需项，并不等于在数值上它必然压倒传统长程分量；
\item 该项专属于 $\Delta L = 2$ 的双核子超弱通道，绝不可将其与单 $\beta$ 衰变中源自介子交换流的两体弱轴流（2BC）有效算符淬灭同名混同。
\end{enumerate}

\section{$\beta$ 缓发粒子发射的马尔可夫吸收网络与多重度归一}\label{sec:appD-delayed-neutron-network}

\subsection{相继级联退激的吸收马尔可夫链模型}

考虑前驱母核经 $\beta$ 跃迁馈送至子核连续谱激发态后的相继退激网络。由于原子核激发能级的退激速率远快于前驱母核的弱衰变周期，处于复合态的核体系退激可严格表述为一个有限状态离散-连续混合的\textbf{吸收马尔可夫链（Absorbing Markov Chain）}模型。

将体系状态空间划分为两类互斥集合：
\begin{itemize}
\item 瞬态子空间（Transient States, $\mathcal{S}_T$）：处于高激发能区、尚未完成全部粒子或辐射退激的所有复合态中间能级；
\item 吸收态子空间（Absorbing States, $\mathcal{S}_A$）：已完全落入束缚基态或超长寿命异能态的终态原子核体系，按照发射中子总数划分为 $x = 0, 1, 2, \dots$ 的稳定终态。
\end{itemize}

马尔可夫转移概率矩阵 $P_{\mathrm{trans}}$ 表达为经典标准分块形式：
\begin{equation}\label{eq:appD-markov-transition-matrix}
P_{\mathrm{trans}} = \begin{pmatrix} \mathbf{Q} & \mathbf{R} \\ \mathbf{0} & \mathbf{I} \end{pmatrix},
\end{equation}
其中 $\mathbf{Q}$ 为瞬态能级之间的单步转移矩阵，其元素由各分支的部分衰变宽度决定：
\begin{equation}\label{eq:appD-q-matrix-element}
Q_{ab} = \frac{\Gamma_{a \to b}}{\Gamma_{\mathrm{tot}}(a)},
\end{equation}
$\mathbf{R}$ 为瞬态向吸收终态跃迁的单步概率矩阵，$\mathbf{I}$ 为吸收态的单位矩阵。

\subsection{终态吸收概率与全多重度归一化证明}

定义吸收马尔可夫链的基本逆矩阵（Fundamental Matrix）为：
\begin{equation}\label{eq:appD-fundamental-matrix}
\mathbf{N} \equiv (\mathbf{I} - \mathbf{Q})^{-1} = \sum_{m=0}^\infty \mathbf{Q}^m.
\end{equation}
矩阵元 $N_{ab}$ 的严格物理意义是：从瞬态 $a$ 出发的级联在被最终吸收之前，平均访问瞬态 $b$ 的期望次数。
从瞬态 $a$ 出发最终被特定的吸收终态（准确发射 $x$ 个中子）所吸收的终态概率 $\Pi_x(a)$ 由矩阵乘积给出：
\begin{equation}\label{eq:appD-absorption-prob-matrix}
\mathbf{B} = \mathbf{N} \mathbf{R} \implies \Pi_x(a) = \sum_b N_{ab} R_{b, x}.
\end{equation}
该吸收概率自然满足马尔可夫平衡差分方程（第~4.7~节式~\eqref{eq:4-7-cascade-pi-def}）：
\begin{equation}\label{eq:appD-pi-balance}
\Pi_x(a) = \delta_{x0} \delta_{at} + \sum_b \frac{\Gamma_{a \to b}}{\Gamma_{\mathrm{tot}}(a)} \Pi_x(b).
\end{equation}
由于体系无无限滞留能级，状态最终必然落入某一个确定的中子发射终态中，因而吸收概率满足严格的条件全概归一性：
\begin{equation}\label{eq:appD-pi-sum-unity}
\sum_{x \ge 0} \Pi_x(a) \equiv 1 \quad (\forall a \in \mathcal{S}_T).
\end{equation}

令 $I_\beta(i)$ 为前驱母核直接弱相互作用馈送至能级 $i$ 的无量纲概率，其母态总分支比满足：
\begin{equation}\label{eq:appD-ibeta-sum}
\sum_{i} I_\beta(i) = b_\beta \le 1.
\end{equation}
按每次母态总衰变归一的最终 $x$ 中子发射净几率为：
\begin{equation}\label{eq:appD-pxn-total-prob}
P_{xn} = \sum_{i} I_\beta(i) \Pi_x(i).
\end{equation}
对所有发射多重度求和：
\begin{equation}\label{eq:appD-pxn-sum-proof}
\sum_{x=0}^\infty P_{xn} = \sum_{x=0}^\infty \sum_i I_\beta(i) \Pi_x(i) = \sum_i I_\beta(i) \left[ \sum_{x=0}^\infty \Pi_x(i) \right] = \sum_i I_\beta(i) \cdot 1 = b_\beta \le 1,
\end{equation}
严格证明了全中子发射多重度概率守恒。

\subsection{阈上竞争与离心势垒阻碍的教学模型}

为了直观证明为何“$Q_\beta > S_n$ 并不等于必定发射中子”，考虑如下简化的教学二能级竞争模型：
设母核衰变布居了子核处于中子分离能微小上方的一个激发态 $|a\rangle$，其能量为 $E_x = S_n + \Delta E$（出射中子可用动能为微弱的 $\Delta E \ll 1\,\text{MeV}$）。
中子发射的部分衰变宽度正比于出射中子的动量与势垒穿透因子 $T_l(\Delta E)$：
\begin{equation}\label{eq:appD-neutron-width-scaling}
\Gamma_n(\Delta E) \propto \sqrt{\Delta E} \cdot T_l(\Delta E).
\end{equation}
若初态自旋与孙核基态之间的自旋选择定则要求出射中子携带较高的轨道角动量（如 $d$ 波 $l=2$ 或 $f$ 波 $l=3$），离心势垒 $V_l(r) = \frac{\hbar^2 l(l+1)}{2\mu r^2}$ 导致穿透几率在低能极限下急剧衰减：
\begin{equation}\label{eq:appD-centrifugal-suppression}
T_l(\Delta E) \propto (\Delta E)^{l + 1/2} \implies \Gamma_n(\Delta E) \propto (\Delta E)^{l + 1} \to 0.
\end{equation}
相反，退激回到子核深束缚态的辐射电磁跃迁宽度正比于光子能量的三次方或五次方（$\Gamma_\gamma \propto E_\gamma^{2\lambda + 1}$），在 MeV 能区具有相对恒定的几率（典型为 $\Gamma_\gamma \sim 0.1$–$1\,\text{eV}$）。
因此，单步中子发射分支比为：
\begin{equation}\label{eq:appD-branch-competition}
\frac{\Gamma_n}{\Gamma_{\mathrm{tot}}} = \frac{\Gamma_n}{\Gamma_n + \Gamma_\gamma} \ll 1.
\end{equation}
此时尽管能态处于连续谱之上，原子核依然以超过 $90\%$ 的几率通过发射级联 $\gamma$ 光子退激回子核束缚态，使得实际单中子发射率远低于简单的能窗总积分。这一教学推演确凿证明了在天体 $r$ 过程核数据反演中必须采用完整的微观统计反应模型（如 Hauser-Feshbach 理论）。

\section{开放环境布居网络、两态平衡与特征本征模}\label{sec:appD-open-network}

\subsection{开放两态动力学主方程}

在天体物理极端环境（如超新星爆发、中子星并合的高温恒星光子热浴）或在束核反应驱动下，原子核体系不断与外界环境发生非弹性散射、光核激发与粒子交换，构成典型的非平衡开放系统。
设体系由两个相互连接的低激发物理状态 $|1\rangle$ 与 $|2\rangle$ 构成（例如同核异能态与基态，能量差为 $\Delta E$）。状态之间的微观布居转移率分别为 $a$（从 $|1\rangle \to |2\rangle$）与 $b$（从 $|2\rangle \to |1\rangle$）。各状态同时通过外部衰变漏出通道（Sinks，如粒子发射、裂变或向外界环境的非可逆退激）具有损耗率 $\Gamma_1$ 与 $\Gamma_2$，并接收外部连续核反应的恒定注入流（Sources）$S_1$ 与 $S_2$。

描述两态粒子数布居 $N_1(t)$ 与 $N_2(t)$ 演化的线性常微分主方程组为：
\begin{equation}\label{eq:appD-master-eq-2state}
\frac{\mathrm{d}}{\mathrm{d}t} \begin{pmatrix} N_1(t) \\ N_2(t) \end{pmatrix} = \mathbf{A} \begin{pmatrix} N_1(t) \\ N_2(t) \end{pmatrix} + \begin{pmatrix} S_1 \\ S_2 \end{pmatrix},
\end{equation}
其中系统耦合演化矩阵为：
\begin{equation}\label{eq:appD-coupling-matrix}
\mathbf{A} = \begin{pmatrix} -(a + \Gamma_1) & b \\ a & -(b + \Gamma_2) \end{pmatrix}.
\end{equation}

\subsection{特征本征值、弛豫模与闭合交换极限}

求解矩阵 $\mathbf{A}$ 的本征值代数方程 $\det(\lambda \mathbf{I} - \mathbf{A}) = 0$：
\begin{equation}\label{eq:appD-char-eq}
\lambda^2 + (a + b + \Gamma_1 + \Gamma_2) \lambda + \left[ (a + \Gamma_1)(b + \Gamma_2) - a b \right] = 0.
\end{equation}
该二次方程的两个特征根为：
\begin{equation}\label{eq:appD-eigenvalues-2state}
\lambda_\pm = -\frac{1}{2}(a + b + \Gamma_1 + \Gamma_2) \pm \frac{1}{2} \sqrt{(a - b + \Gamma_1 - \Gamma_2)^2 + 4 a b}.
\end{equation}
判别式展开为平方和形式：$[(a - b) + (\Gamma_1 - \Gamma_2)]^2 + 4 a b > 0$，因此两个本征根严格为实数且均满足 $\lambda_\pm \le 0$（保证了物理系统在时间上的稳定耗散）。

考察两个极限物理情形：
\begin{enumerate}
\item \textbf{闭合保守交换极限（$\Gamma_1 = \Gamma_2 = 0$ 且无外部注入）：}
此时体系粒子数守恒（$\frac{\mathrm{d}}{\mathrm{d}t}(N_1 + N_2) = 0$）。本征根退化为：
\begin{equation}\label{eq:appD-closed-eigenvalues}
\lambda_1 = 0, \qquad \lambda_2 = -(a + b).
\end{equation}
零本征值 $\lambda_1 = 0$ 对应于总粒子数守恒下的定态渐近分布；而负本征值 $\lambda_2$ 决定了内部粒子交换的弛豫阻尼时间尺度：
\begin{equation}\label{eq:appD-texch-def}
\tau_{\mathrm{exch}} \equiv \frac{1}{a + b}.
\end{equation}
系统任意初始不平衡布居均以特征时间 $\tau_{\mathrm{exch}}$ 指数收敛至由比值 $b/a$ 决定的平衡构型。
\item \textbf{开放耗散分层极限（$a, b \gg \Gamma_1, \Gamma_2$）：}
当内部两态交换速率远大于外部损耗率时，本征值按一阶 Taylor 展开：
\begin{equation}\label{eq:appD-open-fast-slow}
\lambda_{\mathrm{fast}} \approx -(a + b), \qquad \lambda_{\mathrm{slow}} \approx -\frac{b \Gamma_1 + a \Gamma_2}{a + b}.
\end{equation}
体系动力学清晰分裂为两个因果解耦的时间尺度：快模以时间尺度 $\tau_{\mathrm{fast}} = 1/(a+b)$ 迅速使两态达到局部平衡比例 $N_1/N_2 \approx b/a$；随后，慢模以有效加权耗散率 $\Gamma_{\mathrm{eff}} = \frac{b \Gamma_1 + a \Gamma_2}{a + b}$ 控制整个两态耦合包络向外界环境发生长期的不可逆衰变。
\end{enumerate}

\subsection{定态正定性与细致平衡边界}

在存在持续外部注入（$S_1, S_2 > 0$）的定态下，令微分项为零：$\mathbf{A} \mathbf{N}_{\mathrm{ss}} + \mathbf{S} = 0$。由 Cramer 法则，定态解为：
\begin{align}
N_{1, \mathrm{ss}} &= \frac{(b + \Gamma_2) S_1 + b S_2}{(a + \Gamma_1)(b + \Gamma_2) - a b} = \frac{(b + \Gamma_2) S_1 + b S_2}{a \Gamma_2 + b \Gamma_1 + \Gamma_1 \Gamma_2}, \label{eq:appD-n1-ss} \\
N_{2, \mathrm{ss}} &= \frac{a S_1 + (a + \Gamma_1) S_2}{a \Gamma_2 + b \Gamma_1 + \Gamma_1 \Gamma_2}. \label{eq:appD-n2-ss}
\end{align}
由于各项物理速率 $a, b, \Gamma_i, S_i$ 均为非负实数，分母 $a \Gamma_2 + b \Gamma_1 + \Gamma_1 \Gamma_2 > 0$ 恒大于零，保证了定态粒子数在全物理参数空间内的绝对正定性（$N_{i, \mathrm{ss}} \ge 0$）。

在热力学平衡的黑体辐射光子浴中，两态之间的跃迁率满足爱因斯坦自发辐射与受激吸收的微观可逆性，形成著名的\textbf{细致平衡原理（Detailed Balance）}：
\begin{equation}\label{eq:appD-detailed-balance}
\frac{a}{b} = \frac{g_2}{g_1} \exp\left( -\frac{\Delta E}{k_B T} \right),
\end{equation}
其中 $g_1, g_2$ 为态的自旋简并度。
必须澄清该原理的适用边界：细致平衡关系式~\eqref{eq:appD-detailed-balance} 严格依赖于光子浴处于完全热动平衡状态。在低能核反应在束实验、高能激光等离子体碰撞或中微子非热退耦环境下，外场具备强相干性或定向流，细致平衡被强烈打破；此时两态演化必须由完整的开放量子主方程或 Lindblad 耗散子统一求解。
